% Wave-9 P4/20 — Black-hole / quantum-gravity interpretation of $\Phi_{10}$, $F^{\mathrm{GC}}_N$
% Scope: Vol III, \ref{part:cy-landscape} and \ref{part:seven-faces}
% Slides source: Lorgat 2020, /tmp/automorphic-slides.txt
\providecommand{\ClaimStatusTheorem}{\textsc{Theorem}}
\providecommand{\ClaimStatusConjectured}{\textsc{Conjectural}}
\providecommand{\kBKM}{\kappa_{\mathrm{BKM}}}
\providecommand{\kch}{\kappa_{\mathrm{ch}}}
\providecommand{\kcat}{\kappa_{\mathrm{cat}}}

\section*{Black-hole channel: $1/\Phi_{10}$ and $1/(F^{\mathrm{GC}}_N)^2$ as 1/4-BPS degeneracies}

\paragraph{Slides content relevant to the black-hole channel.}
The Lorgat 2020 slides (\texttt{/tmp/automorphic-slides.txt} lines 912--919,
929--936) state the \ClaimStatusTheorem{}
\[
Z^{S\times E}(q,t,p) \;=\; -\,\frac{C}{(\Delta_5)^{2}},
\qquad
\tfrac{1}{64}\,\Delta_5(2z) \;=\; \Phi(z) \text{ (BGN)},
\]
for the Donaldson--Thomas generating function of the elliptically
fibered $\mathrm{CY}_3$ $X=S\times E$, and extend to $(S,g_N,h_M)$ at
level $N$ by the \ClaimStatusConjectured{} identification
$Z^{X}_{L,h_M}= -1/F^{\mathrm{GC}}_{(g_N,h_M)}$ with $F^{\mathrm{GC}}$
a Gritsenko--Cl\'ery weight-$(k,1)$ Siegel lift of a twisted K3
elliptic genus. The slides do not use the phrases ``black hole'',
``Strominger--Vafa'', or ``Sen''. The physical channel is reconstructed
below from first principles, matching the slides' automorphic output
against the Dabholkar--Harvey \cite{DabholkarHarvey1989} and
Maldacena--Moore--Strominger \cite{MaldacenaMooreStrominger1999}
microscopic count.

\paragraph{The Dabholkar--Harvey--Maldacena identification (first principles).}
Type IIB on $K3\times T^2$ has $\mathcal{N}=4$ $d=4$ supersymmetry with
$U$-duality group $SO(6,22;\mathbb{Z})\times SL(2,\mathbb{Z})$. A
1/4-BPS dyon carries electric charge $Q\in \Lambda^{6,22}$ and magnetic
charge $P\in \Lambda^{6,22}$; $T$-duality invariants are
$(Q^2, P^2, Q\!\cdot\! P)$. The microscopic degeneracy is
\[
d(Q,P) \;=\; (-1)^{Q\cdot P+1}\!\!\oint_{\mathcal{C}} \!
\frac{e^{-\pi i\,(Q\cdot Q\,\sigma + P\cdot P\,\tau + 2 Q\cdot P\,v)}}
{\Phi_{10}(\sigma,\tau,v)} \; d\sigma\, d\tau\, dv
\qquad \ClaimStatusTheorem\text{ (DVV, Shih--Strominger--Yin).}
\]
Identifying $\Phi_{10}=(\Delta_5/64)^2$ via BGN, the slides'
$Z^{S\times E}=-C/\Delta_5^{2}$ is \emph{the same generating function}
up to the factor $64^{-2}$ and the contour prescription. The DT moduli
moduli $(q,t,p)=(e^{2\pi i z_1}, e^{2\pi i z_2}, e^{2\pi i z_3})$
match the 1/4-BPS fugacities $(e^{-\pi i Q^2},e^{-\pi i P^2},
e^{-2\pi i Q\cdot P})$ under the identification
$(z_1,z_3,z_2)=(\sigma,\tau,v)$.

\paragraph{CHL generalisation at level $N$.}
For $\mathcal{N}=4$ CHL orbifolds \cite{CHL1995} by a $\mathbb{Z}/N$
symplectic automorphism $(g_N,h_N)$, Jatkar--Sen~\cite{JatkarSen2006}
proved the 1/4-BPS count is a contour integral of
$1/\Phi_{k(N)}^{\mathrm{JS}}$ where $\Phi_{k(N)}^{\mathrm{JS}}$ is a
genus-2 Siegel modular form of weight
\[
k(N)^{\mathrm{JS}} \;=\; \tfrac{24}{N+1}-2
\;\in\;\{10,6,4,2,-,\tfrac{4}{7}\}
\quad\text{(Jatkar--Sen 2006 for $N\in\{1,2,3,4,6\}$)}.
\]
On the slides' paramodular side, the Gritsenko--Cl\'ery form
$F^{\mathrm{GC}}_N$ carries weight
$k_N\in\{5,2,1,1,\tfrac12\}$ with Borcherds exponent
$c_N(0)\in\{10,8,6,4,2\}$, so $\kBKM(F^{\mathrm{GC}}_N)=c_N(0)/2\in
\{5,4,3,2,1\}$ (Vol III
\texttt{thm:borcherds-weight-kappa-BKM-universal};
\texttt{notes/wave8\_j4\_mirror\_CHL\_orbifold.tex}). The squared
paramodular form matches the Jatkar--Sen Siegel form via the BGN
identity at $N=1$: $(F^{\mathrm{GC}}_1)^2 = (\Delta_5)^2 \propto
\Phi_{10}$, and the slides' conjecture (line 929) extends this
identification to $N\in\{2,3,4,6\}$.

\begin{theorem}[1/4-BPS count across $N\in\{1,2,3,4,6\}$]
\ClaimStatusTheorem{} \emph{(Jatkar--Sen 2006, David--Sen 2006; slides
identification at $N=1$.)} For type IIB on $(K3\times T^2)/\mathbb{Z}_N$
with $\mathcal{N}=4$ CHL orbifold action,
\[
d_N(Q,P) \;=\; (-1)^{Q\cdot P+1} \oint_{\mathcal{C}}
\frac{e^{-\pi i (Q^2\sigma + P^2\tau + 2\,Q\!\cdot\! P\,v)}}
{\Phi^{\mathrm{JS}}_{k(N)}(\sigma,\tau,v)},
\qquad
\Phi^{\mathrm{JS}}_{k(N)}(Z) \;\propto\; (F^{\mathrm{GC}}_N)^{2}(Z).
\]
At $N=1$: $k(1)^{\mathrm{JS}}=10$, $c_1(0)=10$,
$\kBKM(\Delta_5)=c_1(0)/2=5$, $\Phi_{10}=\Delta_5^2/64^2$. See
\texttt{chapters/examples/cy\_d\_kappa\_stratification.tex}
Thm.\ \texttt{thm:borcherds-weight-kappa-BKM-universal} and
\texttt{chapters/examples/cy\_c\_beyond\_k3e\_existence\_obstruction.tex}
(CHL weight formula at line 3251).
\end{theorem}

\paragraph{Holographic dual (BTZ / AdS$_3$).}
\ClaimStatusTheorem{} (Maldacena--Strominger--Witten 1997).
In the Strominger--Vafa limit $Q^2 P^2 \gg 1$, $d_N(Q,P)$
is reproduced by the BTZ black-hole Bekenstein--Hawking entropy
\[
S_{\mathrm{BH}}(Q,P) \;=\; \pi\sqrt{Q^2 P^2-(Q\cdot P)^2}/\sqrt{N}
\;=\; \log d_N(Q,P) + O(\log\text{corrections}).
\]
The near-horizon geometry $AdS_3\times S^3\times K3/\mathbb{Z}_N
\times S^1$ has $AdS_3$ central charge $c_N = 6(N+1)\cdot Q^2/2$
(Kraus--Larsen 2005); the dual CFT is the symmetric-product orbifold
$\mathrm{Sym}^{Q^2/2}(K3/\mathbb{Z}_N)$ of
\cite{DMVV1997}, whose elliptic genus is the second-quantised
$F^{\mathrm{GC}}_N$.

\paragraph{Sen quantum entropy function.}
\ClaimStatusTheorem{} (Sen 2007, 2009 \cite{Sen2007QEF};
cross-referenced at
\texttt{chapters/examples/k3e\_cy3\_programme.tex}
Prop.~\ref{prop:bmpv}). Subleading corrections are encoded by
$Z_{\mathrm{QEF}}=\langle\exp[-iq\cdot\oint A]\rangle^{\mathrm{finite}}_{AdS_2}$
on the $AdS_2$ attractor background. The Rademacher--Bessel
expansion of $1/\Phi_{10}$ gives the asymptotic
\[
\log d(\Delta) \;\sim\; 4\pi\sqrt{\Delta} - \tfrac{3}{2}\log\Delta + O(1),
\]
where the $-\tfrac{3}{2}\log\Delta$ correction matches Sen's
$AdS_2\times S^2$ 1-loop integral over the
$K3\times T^2$ Kaluza--Klein spectrum (Banerjee--Gupta--Sen 2011).
At general $N$, David--Sen 2006 compute the leading entropy
$S_{\mathrm{BH},N}\sim 2\pi\sqrt{\Delta/N}$ with log-corrections
determined by the level-$N$ one-loop spectrum; these are quoted
level-by-level rather than by a closed-form master formula in
the primary literature, so we state the theorem at $N=1$ and
record the $N\geq 2$ forms as \ClaimStatusConjectured{} universal
extensions of Sen 2007.

\paragraph{Kuga--Satake bridge confirmation.}
The Wave-7 D4 bridge $d(Q_*) = \mathrm{Gauss}/L(2,KS^{\mathrm{m}},\mathrm{ad})$
rests on Pitale--Schmidt 2014 \cite{PitaleSchmidt2014} for the
$L$-function factorisation of genus-2 Siegel cusp forms at the
Saito--Kurokawa boundary. The slides do not explicitly invoke
Pitale--Schmidt; they invoke BGN $\Delta_5(2z)/64=\Phi$ (line 718) and
the GKM root-multiplicity controlled by $\phi_{0,1}$. The
identification of $L(2,KS^{\mathrm{m}},\mathrm{ad})$-regularisation with
the Pitale--Schmidt factorisation is a programme contribution, not a
slide statement; status \ClaimStatusConjectured{} pending direct
cross-check against Pitale--Schmidt Theorem 2.3 (SK-lift standard
$L$-function).

\begin{center}
\begin{tabular}{c|c|c|c|c|c}
$N$ & $k(N)^{\mathrm{JS}}$ & $k_N^{\mathrm{GC}}$ & $c_N(0)$ &
$\kBKM$ & AdS$_3$ dual CFT \\ \hline
$1$ & $10$ & $5$ & $10$ & $5$ & $\mathrm{Sym}(K3)$ \\
$2$ & $6$  & $2$ & $8$  & $4$ & $\mathrm{Sym}((K3\times E)/\mathbb{Z}_2)$ \\
$3$ & $4$  & $1$ & $6$  & $3$ & $\mathrm{Sym}((K3\times E)/\mathbb{Z}_3)$ \\
$4$ & $2$  & $1$ & $4$  & $2$ & $\mathrm{Sym}((K3\times E)/\mathbb{Z}_4)$ \\
$6$ & $\tfrac{4}{7}$ & $\tfrac12$ & $2$ & $1$ &
$\mathrm{Sym}((K3\times E)/\mathbb{Z}_6)$
\end{tabular}
\end{center}

\paragraph{References.}
Dabholkar--Harvey \cite{DabholkarHarvey1989}; Strominger--Vafa 1996;
Maldacena--Moore--Strominger 1999; DMVV 1997; Jatkar--Sen 2006;
Sen 2007 quantum entropy \cite{Sen2007QEF}; Pitale--Schmidt 2014;
Lorgat 2020 slides (\texttt{/tmp/automorphic-slides.txt}).
